% --- CAPÍTULO 1 ---
\chapter{Introducción a la gestión de Recursos Humanos}

La función de personal ha evolucionado de ser una tarea de planificación y control a convertirse en una función estratégica de integración, coordinación y apoyo. Este cambio implica comprender las expectativas y sentimientos del personal, pasando de la antigua \textbf{concepción taylorista} (estímulos puramente económicos) a considerar al ser humano como un \textbf{factor estratégico} indispensable.

Aunque la función ha existido siempre, el interés formal por su gestión surge en el siglo XX debido al aumento en la complejidad de las organizaciones.

\section{Factores que influencian la gestión de Recursos Humanos}

La gestión de RR. HH. está condicionada por tres dimensiones fundamentales:

\subsection{Factores del entorno}
Son fuerzas externas fuera del control de la dirección que generan amenazas y oportunidades. Las empresas deben ser flexibles para reaccionar ante ellos:

\begin{itemize}
    \item \textbf{Rapidez de los cambios:} El entorno actual es global, dinámico y complejo.
    \item \textbf{Crecimiento de Internet:} Ha transformado la comunicación y creado nuevas oportunidades de negocio.
    \item \textbf{Legislación:} Obligación de cumplir con la normativa laboral para evitar sanciones.
    \item \textbf{Conciliación de la vida laboral y familiar:} Políticas que permiten equilibrar responsabilidades.
    \item \textbf{Diversidad de la fuerza de trabajo:} Gestión de diferentes perfiles para aprovechar nuevas habilidades.
\end{itemize}

\subsection{Factores organizativos}
Elementos internos que pueden ser controlados por la dirección:

\begin{itemize}
    \item \textbf{Posición competitiva:} Atraer talento para mejorar la productividad e innovación.
    \item \textbf{Flexibilidad empresarial:} Lograda mediante descentralización o nuevas tecnologías.
    \item \textbf{Reestructuración organizativa:} Cambios en la estructura para mejorar la eficiencia.
    \item \textbf{Tamaño empresarial:} A mayor tamaño, más recursos pero mayores retos de coordinación.
    \item \textbf{Cultura organizativa:} Valores y normas compartidas; si es rígida, puede limitar la innovación.
\end{itemize}

\subsection{Factores individuales}
Dependen de las características personales de los empleados:

\begin{itemize}
    \item \textbf{Ajuste Persona-Organización:} Compatibilidad entre el empleado y la cultura de la empresa.
    \item \textbf{Ética y responsabilidad social:} Afecta a empleados y otros grupos de interés.
    \item \textbf{Productividad del empleado:} Resultado de sus capacidades, motivación y condiciones.
    \item \textbf{Empowerment (Delegación de poder):} Otorgar capacidad de decisión para aumentar el compromiso.
    \item \textbf{Fuga de cerebros:} Emigración de profesionales altamente capacitados.
    \item \textbf{Inseguridad laboral:} La percepción de riesgo en el empleo afecta el bienestar y rendimiento.
\end{itemize}



\section{La gestión de los recursos humanos}
Es el proceso de atraer, desarrollar, motivar y retener a los trabajadores.

\subsection{Objetivos de la gestión de RR. HH.}
\begin{enumerate}
    \item \textbf{Objetivos explícitos:} Atraer candidatos, retener empleados valiosos, motivar el compromiso y desarrollar habilidades.
    \item \textbf{Objetivos implícitos:} Mejorar la productividad, asegurar la calidad de vida laboral y cumplir la normativa (seguridad, higiene, pagos).
    \item \textbf{Objetivos a largo plazo:} Rentabilidad, competitividad, incremento del valor de la empresa y eficiencia organizativa.
\end{enumerate}

\subsection{Procesos de la gestión de RR. HH.}
Los procesos se dividen en cinco bloques principales:
\begin{enumerate}
    \item \textbf{Básicos:} Análisis del puesto y planificación.
    \item \textbf{Afectación:} Reclutamiento, selección y socialización.
    \item \textbf{Formación y desarrollo:} Formación y gestión de la carrera profesional.
    \item \textbf{Evaluación y compensación:} Evaluación del desempeño y sistemas retributivos equitativos.
    \item \textbf{Sustractivos:} Despidos, dimisiones y jubilaciones.
\end{enumerate}

\subsection{El departamento y el director de recursos humanos}
El departamento es el área encargada de la ejecución de estos procesos, liderada por un director responsable de asegurar el cumplimiento de los objetivos estratégicos.

\section{Tendencias actuales}
\begin{itemize}
    \item \textbf{Gestión de la diversidad:} Estrategia integral para atraer perfiles diferentes que aporten nuevos puntos de vista.
    \item \textbf{Gestión del talento:} Procesos para atraer y retener empleados con el fin de cubrir futuros puestos directivos.
\end{itemize}